# Profil J2 — Prise de notes brute (Analyste, 6 mois)

**Code :** J2  
**Fonction :** Analyste  
**Ancienneté :** 6 mois  
**Contexte :** Gestion de crise, appels d'offre, institutions européennes.

---

## T0 — Cadrage

**Q : Parcours, comment arrivé ici ?**
- 6 mois. Avant, stage armée. Gestions de crise. Aimais situations incertitude.
- Propales. PPT de réponse. Préparation exos (PPT, Word). Jour J : déroulé. Après : analyse Word/Excel, restitution PPT.
- Cycle complet.

---

## T1 — Usage réel

**Q : Dernière fois utilisé outil génératif ?**
- Propale. 3 jours pour répondre. Section "leviers d'innovation".
- Copilot PPT. Fils contexte 2 phrases. Demandé 5 pistes d'injects innovants.
- Sorti "panne coms", "fuite presse". Certaines utilisables, d'autres trop génériques.
- Gardé 3. Adapté formulations. Construit slide. Gagné temps brainstorming.

**Q : Sur quelles tâches au quotidien ?**
- Propales. Injects. Créatif : scénarios, rebondissements, données fictives crédibles.
- Jamais sur analyse post-exercice. Chiffrage, REX concrets. L'IA n'était pas dans la salle. Fais main.

**Q : Comment appris ?**
- Sur le tas. Pas de formation. Intuitif. Tu demandes, il répond. Pas de prompts avancés. Simple.

**Q : Temps gagné mis sur quoi ?**
- Gagné 1h sur recherche d'idées. Utilisée pour affiner structure, vérifier arguments. Passé + temps sur cohérence globale. Bon équilibre.

**Q : Situation où plus d'effort que de le faire soi-même ?**
- Oui. Injects cyber hyper spécifique. Trop générique. DDoS, ransomwares...
- 10min à reformuler prompt, exemples. Rien. Réécrit en 5min. Trop spécifique = tête.

**Q : Vous en parlez entre vous ?**
- Un peu. Astuces. Informel. Pas en réunion ou avec managers.
- Tout le monde l'utilise, personne en parle. Secret de polichinelle.

---

## T2 — Apprentissage

**Q : Comment on apprend le métier aujourd'hui ?**
- Pratique. Propales, exos, restitutions. IA aide à aller vite, mais n'apprend pas le métier.
- Clients, seniors, analyse terrain. Outil, pas mentor.

**Q : Par quelles tâches on apprend vraiment ?**
- Incertitude. Propale, tu sais pas ce que client attend. Exo, paramètres changent. Adaptation.
- Répétitives = rien. IA fait bien.
- Cœur = décision sous pression. IA le fait pas.

**Q : Tâches ingrates, injects ?**
- Injects = créatif. IA partenaire brainstorming. Sort pistes, évalues, améliores.
- Gagne temps sur génération. Consacre sur affinage/crédibilité. Pas appauvrit, rend + efficace. Mais garde la main.

**Q : Comment vérifiez ce que produit l'outil ?**
- Cohérence globale. Crédibilité dans contexte client. Si incohérences ou trop générique, écarte ou retravaille.
- Pas vérifier sources/chiffres. Pas factuel. Créativité/structure.

---

## T3 — Client et valeur

**Q : Confiance client repose sur quoi ?**
- Anticiper. Scénarios, imprévus, conséquences.
- Plan solide + adaptation en direct = confiance.
- IA prépare, jour J consultant terrain.

**Q : Diriez-vous au client qu'une partie est IA ?**
- Non. N'apporte rien. Client veut livrable de qualité. Peu importe comment.
- Si dis "utilisé Copilot", distraction. "C'est vraiment vous ?". Pas mentionner.

**Q : Livrable moins crédible ?**
- Ancrage. Si client sent pas compris secteur, contraintes, culture, doute.
- IA générique, pas nuances contexte. Livrable juste mais pas pertinent. Garde main sur analyse/jugement.

**Q : Dans 10 ans, client paiera pour quoi ?**
- Expérience humaine. Crise = relationnel, improvisation, prise de risque.
- IA prépare, simule. Pas dans salle. Pas sentir tension, lire visages, ajuster.
- Conseil = présence, pas algorithmes.

---

## T4 — Gouvernance

**Q : Règles sur usage ?**
- Copilot officiel. Communication générale. Rien contraignant.
- Passé armée => sensibilisé confi. Pas besoin qu'on dise. Je fais nature.

**Q : Avant livrable client ?**
- Relecture manager/senior. Cohérence, chiffres, reco.
- Pas demandent si IA. S'en doutent peut-être. Contenu bon, processus pas regardé.

**Q : Déjà rencontré contenu faux IA ?**
- Pas vraiment. Trop générique, mais pas inventé. Utilise pour idées, pas faits. Risque limité.
- Collègues ChatGPT données = chiffres nulle part. Ont dû reprendre. Je préfère éviter.

**Q : Règle idéale ?**
- "Aide créativité/structure, pas substitut jugement."
- Vérification : si faits, justifiables.
- Données sensibles : pas outils non-officiels. Pas compliqué.

---

## T5 — Clôture

**Q : À la place de la direction, première décision ?**
- Charte simple. 1 page. 3 règles : possible, impossible, vérifier.
- Canal Teams bonnes pratiques. Actuellement on tâtonne. Certains font erreurs évitables.

**Q : Aspect important non abordé ?**
- Culture cabinet. Armée = IA mal vue. Ici ouvert, mais ambivalence. On dit innover, pas dit comment intégrer. Décalage discours/pratique.

**Q : Point de vue utile à recueillir ?**
- Consultant 10 ans gestion crise. Sait pas comment il perçoit outils. Il les utilise ? Craint ? Voit différence juniors IA vs non-IA ? Regard éclairant.